\subsection{Sample Validation}\label{sec:verification}

Nugget addresses the limitations of using simulation to validate selected samples (Section~\ref{sec:validating_the_selected_samples_is_infeasible}) by generating \emph{nuggets} that can run directly on both real hardware and simulation.
Therefore, we can validate the selected samples in real hardware, and use them for simulation.
This approach allows for fast, scalable validation of sampling accuracy without the high overhead of simulation, even for large applications.

Because nuggets are \textbf{binary-independent}, they can be validated across different systems with minimal constraints. 
Validating the same samples on multiple platforms enhances confidence in their \emph{representativeness and robustness}. 
As shown later in Section~\ref{sec:model_performance_evaluation}, we found that having \emph{consistent prediction error across platforms} is a stronger indicator of sample quality than achieving low error on a single platform.

To run a nugget on real hardware, the system must still execute the portion of the program leading up to the sample's start marker (Section~\ref{sec:markers}).
For samples that are in later stages of a large program's execution, this can introduce some validation overhead. 

In the future, we plan to reduce this overhead further by exploring techniques such as creating software checkpoints using CRIU~\cite{CRIU} \revision{or BLCR~\cite{BLCR}} to resume execution closer to the sample start point.
